% Generated draft: requires mathematical and source review.
\documentclass[11pt]{article}
\usepackage[margin=25mm]{geometry}
\usepackage{fontspec}
\setmainfont{Latin Modern Roman}
\usepackage{amsmath,amssymb,mathtools,unicode-math}
\setmathfont{Latin Modern Math}
\usepackage{enumitem,needspace,xurl,hyperref}
\setlength{\parindent}{0pt}
\setlength{\parskip}{5pt}
\setlength{\emergencystretch}{3em}
\newcommand{\N}{\mathbb{N}}
\newcommand{\Z}{\mathbb{Z}}
\newcommand{\Q}{\mathbb{Q}}
\newcommand{\R}{\mathbb{R}}
\newcommand{\C}{\mathbb{C}}
\newcommand{\F}{\mathbb{F}}
\newcommand{\A}{\mathbb{A}}
\newcommand{\PP}{\mathbb{P}}
\newcommand{\E}{\mathbb{E}}
\newcommand{\eps}{\varepsilon}
\newcommand{\dd}{\,\mathrm{d}}
\begin{document}
\section*{UnsolvedMath ID: 30000752}
\textbf{Source problem code:} catalogue entry
\subsection{Definitions and mathematical statement}
For a complex parameter $q$ with $|q|>1$, define the Tschakaloff function
\[
T_q(z)=\sum_{\nu=0}^{\infty}q^{-\nu(\nu+1)/2}z^\nu.
\]
The series converges absolutely for every $z\in\C$. Fix the rational argument $z$ intended by the source, and let $\mathcal A_z\subseteq\{q\in\C:|q|>1\}$ denote its intended set of admissible parameters. The proposed linear independence over $\Q$ means
\[
\begin{gathered}
\text{for every finite list of distinct }q_1,\ldots,q_s\in\mathcal A_z
\text{ and every }a_1,\ldots,a_s\in\Q,\\
\sum_{i=1}^{s}a_iT_{q_i}(z)=0\quad\Longrightarrow\quad a_1=\cdots=a_s=0.
\end{gathered}
\]
The argument $z$ is common to these values; the independence question varies $q$. The number $1$ is not adjoined to the list of values, since the source does not ask for that stronger independence statement.

The qualifier ``admissible'' is essential but undefined: the record does not say whether $q$ is integral, rational or algebraic, or state restrictions on parameter relations or on $z$. In particular, $T_q(0)=1$ for every $q$, so the unrestricted assertion for rational $z=0$ and two distinct admissible parameters fails. A nonzero-argument condition or another admissibility convention is required, but is not silently supplied. The source's canonical-version claim does not itself recover these missing conditions.
\subsection{Short English statement}
For a fixed rational argument, prove rational linear independence of Tschakaloff values at distinct admissible parameters. The admissible parameter set is undefined, and the unrestricted zero-argument case gives identical values, so the intended qualifications must be recovered.
\subsection{Sources}
\begin{itemize}
\item Supplied report: \emph{Diophantische Approximationen} (2008), Oberwolfach Reports, DOI 10.4171/owr/2007/21, pp. 1115--1190; proposed year 2007. \url{https://doi.org/10.4171/owr/2007/21}.
\item Supplied Vaananen--Zudilin references: \emph{Russian Mathematical Surveys}, \url{https://doi.org/10.1070/RM2007v062n01ABEH004386}; and \emph{Journal of Number Theory}, \url{https://doi.org/10.1016/j.jnt.2008.03.003}.
\item Local UnsolvedMath registry, permanent ID 30000752, catalogue entry; created 25 August 2026, updated 16 September 2026. Its corrected\_verified note calls this the complete canonical version, including the series, of duplicate OWR-1537-001. That companion is permanent ID 30000749 and is labelled unrecoverable only as the less complete duplicate.
\item The supplied 21 August 2026 triage labels the question open, reports restrictive parameter configurations in existing results, and refers to a 2025 Vaananen--Zudilin note still describing the cross-parameter problem as open. No separate bibliographic identification of that 2025 note is provided beyond the two references above; the dated claim is retained without independent verification.
\end{itemize}
\end{document}
